Magnetic resonance imaging (MRI) is central to both the diagnosis and the
monitoring of multiple sclerosis (MS). Under the 2024 revision of the
McDonald criteria, dissemination in space is established by T2-hyperintense
lesions in the periventricular, juxtacortical or cortical, and
infratentorial regions, the spinal cord, and the optic nerve, and
dissemination in time can be shown by a single new T2-hyperintense lesion on
follow-up MRI~\cite{montalban_diagnosis_2025,barkhof_2024_2025}. After
diagnosis, new or enlarging lesions on follow-up scans are the main
radiological marker of disease activity and inform treatment
decisions~\cite{wattjes_2021_2021}. In both settings, a single lesion can
change the conclusion of a report.

Many of these lesions are small. Lesion size in MS is strongly skewed
towards small lesions~\cite{wang_survey_1997}, and acquisitions that miss
lesions underestimate lesion load~\cite{filippi_quantitative_1996}.
The 2024 MAGNIMS--CMSC--NAIMS recommendations state that lesions should
generally be at least 3~mm in longest diameter, while smaller lesions in
specific locations, such as the corpus callosum or the floor of the fourth
ventricle, can also be considered demyelinating~\cite{barkhof_2024_2025}.
The central vein sign, which the 2024 criteria introduce as a marker of
specificity, is likewise assessed in lesions of at least this
size~\cite{montalban_diagnosis_2025,barkhof_2024_2025}. Lesions just above
3~mm are therefore clinically relevant but are also the hardest to find by
eye: on a whole-brain three-dimensional fluid-attenuated inversion recovery
(FLAIR) scan they occupy a few voxels among millions, and finding them,
in particular new ones on serial scans, depends on careful visual
inspection of the whole volume.

Automated lesion segmentation is designed to support this task. The 2024
recommendations note that automated tools that report the number,
location, and volume of (new) lesions can potentially improve the accuracy
and objectivity of MRI assessment, while their clinical benefit still
requires validation~\cite{barkhof_2024_2025}. Deep learning methods, in
particular nnU-Net~\cite{isensee_nnu-net_2021} and MS-specific tools
such as LST-AI~\cite{wiltgen_lst-ai_2024} and
FLAMeS~\cite{dereskewicz_flames_2025}, now achieve high agreement with
expert annotations on standard overlap metrics. Their errors, however, are
concentrated in small lesions: in FLAMeS, most missed lesions were smaller
than 10~mm$^3$~\cite{dereskewicz_flames_2025}. For a tool that works in
tandem with a human reader, these two kinds of error are not equivalent. A
false-positive detection is shown to the reader and can be dismissed in
seconds, whereas a lesion that the tool does not mark gives the reader no
prompt to look at that location. The most useful property of such a tool
is therefore that it misses as few lesions as possible, with false
positives kept at a level that a reader can review.

A main reason why current models miss small lesions lies in how they are
trained. The standard training objective, which combines the Dice score
with voxel-wise cross-entropy, rewards correctly classified voxels. A
single large lesion contains as many voxels as dozens of small ones, so
during training, a missed 10-voxel lesion costs the model almost nothing,
and a model can reach a high Dice score while leaving many small lesions
undetected. Existing alternatives reweight voxels by class, difficulty, or
position, or balance lesions of different sizes, but they still judge each
lesion by how well it is outlined rather than by whether it is found at
all (\cref{sec:related}).

In this study, we change the training objective rather than the network.
The proposed objective, CATMIL, adds to the standard nnU-Net loss a
lesion-level \emph{Multiple Instance Learning (MIL)} term that penalizes
any lesion in which no voxel is detected, mirroring the question a reader
asks: was this lesion found? It also includes a \emph{Component-Adaptive
Tversky (CAT)} term that gives every lesion a comparable weight regardless
of its volume, so that detected small lesions are also delineated; an
ablation shows that the MIL term drives the detection gain and that
adding the CAT term gives the best voxel-level and boundary accuracy.
Because only the objective changes, CATMIL can be
used in existing nnU-Net-based pipelines without changes to the
architecture or to inference. We evaluated CATMIL on two public MS
datasets with T1-weighted images (T1-WI), T2-weighted images (T2-WI),
and FLAIR images, MSLesSeg and
3D-MR-MS, and report outcomes that correspond to clinical reading:
detection of lesions of at least 3~mm, the number of completely missed
lesions per scan, and the number of false-positive detections.

The main findings of this study are as follows:
\begin{itemize}
\item On MSLesSeg, CATMIL reduces the number of completely missed lesions
by about half compared with the standard nnU-Net objective, in every test
patient and without loss of voxel-level agreement with the reference
segmentation; the reduction is reproduced on the 47 cross-validation
patients and against four alternative training objectives. CATMIL also
improves the detection of lesions of at least 3~mm on both datasets. The
detection gain is due to the lesion-level (MIL) term, while adding the CAT
term gives the best voxel-level and boundary accuracy of the ablation
variants.
\item Standard objectives do not miss small lesions by a narrow margin:
for most of them, the model produces no response at all. Making an
existing tool more sensitive by lowering its decision threshold therefore
cannot recover these lesions; the training objective must change.
\item The cost of the higher sensitivity is a larger number of small
false-positive detections. A simple minimum-size filter removes most of
them while preserving the gain in detected lesions, so the operating point
can be adapted to the intended use, such as a sensitive setting for
follow-up reading with radiologist confirmation. This cost depends on the
dataset: it was considerably larger on 3D-MR-MS, where a larger filter was
needed.
\end{itemize}
